\documentclass[11pt]{article}
\usepackage[a4paper,margin=25mm]{geometry}
\usepackage{amsmath,amssymb,amsthm,mathtools,booktabs,array}
\usepackage[T1]{fontenc}
\usepackage{lmodern,microtype}
\usepackage[hidelinks]{hyperref}
\usepackage{enumitem}
\hypersetup{pdftitle={Uniform eventual decrease of permanents under matrix powers},pdfauthor={Research manuscript generated in this conversation},pdfsubject={Unreviewed proof of Foregger's 1978 per-k conjecture},pdfkeywords={permanent, doubly stochastic matrices, Foregger conjecture}}
\setlist{nosep}
\setlength{\parskip}{4pt}
\setlength{\parindent}{0pt}
\newtheorem{theorem}{Theorem}[section]
\newtheorem{lemma}[theorem]{Lemma}
\newtheorem{corollary}[theorem]{Corollary}
\theoremstyle{remark}\newtheorem{remark}[theorem]{Remark}
\DeclareMathOperator{\per}{per}
\DeclareMathOperator{\dist}{dist}
\newcommand{\Om}{\Omega_n}
\newcommand{\Fset}{\mathcal F_n}
\newcommand{\fn}[1]{\left\|#1\right\|_{\mathrm F}}
\newcommand{\op}[1]{\left\|#1\right\|_{2}}
\newcommand{\one}{\mathbf 1}
\newcommand{\eps}{\varepsilon}
\title{Uniform eventual decrease of permanents\\
under matrix powers\\[5pt]\large A proof of Foregger's 1978 per-$k$ conjecture}
\author{Research manuscript generated in this conversation}
\date{7 September 2026\quad Version 1.0}
\begin{document}
\maketitle
\begin{abstract}
For every $n$, there is an integer $N(n)\ge2$ such that
$\per(A^k)\le\per A$ for every nonnegative doubly stochastic $n\times n$
matrix $A$ and every integer $k\ge N(n)$. Equality, for these exponents,
holds precisely for uniform block-permutation matrices defined below.
This proves the uniform-exponent conjecture attributed to Foregger in
Minc's 1978 catalogue, and strengthens it to all sufficiently large
exponents. The proof combines a categorical-occupancy bound for
row-averaged matrices with a local estimate that remains valid for
arbitrarily large powers. A compactness argument then handles the full
Birkhoff polytope without a positive lower bound on its nonzero entries.
A uniform quadratic stability consequence is also obtained.
The global exponent is not made explicit.
\end{abstract}
\begin{center}\small
\textbf{Review status.} This is an unreviewed mathematical manuscript, not
an accepted publication or an externally certified priority claim.
The general theorem is proved analytically; the accompanying exact
computations are independent implementation-level diagnostics, not a
formal-machine proof of the theorem.
\end{center}

\section{Target, prior results, and precise claim}
Let $\Om$ denote the compact set of real nonnegative matrices with every
row and every column summing to one. All powers in this manuscript are
\emph{ordinary matrix powers}. Write
\[
 f(A)=\per A=\sum_{\sigma\in S_n}\prod_{i=1}^n a_{i,\sigma(i)}.
\]
Foregger's conjecture asks whether, for each $n$, there exists an integer
$k(n)>1$, independent of $A$, for which
\begin{equation}\label{eq:conjecture}
 f(A^{k(n)})\le f(A)\qquad\text{for every }A\in\Om.
\end{equation}
The conjecture appears as Conjecture 17 in Minc's catalogue, and is
restated in Cheon--Wanless \cite{CW} and Zhang \cite{Zhang}. Melnykova
\cite[Section 4.1]{Mel} reports Foregger's original letter to Minc as
dated 17 December 1977; the published catalogue is from 1978. The original
letter and the original book page were not independently inspected here.
The restriction $k>1$, made explicit in \cite{Mel}, is essential.

Chang proved the $n=3$ case, with exponent eight; see \cite{Chang90} and
the account in \cite{Mel}. Melnykova's Theorem 4.10 gives a common
sufficiently-large exponent for matrices whose \emph{nonzero} entries
are bounded below by a fixed $c>0$. Its threshold depends on $n$ and $c$.
Her Corollary 4.11 gives a matrix-dependent exponent for every individual
matrix. Removing dependence on a lower-entry bound or on $A$ is the
issue addressed here. The present proof does not use either partial
result as an unproved uniform assertion. In a dated primary discussion
on 15 May 2022, Wanless again identified this general question as open
\cite{MO}; that statement is not treated as exhaustive 2026 status evidence.

\subsection{The exceptional matrices}
Partition $[n]=C_1\sqcup\cdots\sqcup C_s$ into nonempty classes, put
$m_a=|C_a|$, and let $\pi$ be a permutation of the classes satisfying
$m_{\pi(a)}=m_a$. Define
\begin{equation}\label{eq:Fdef}
 F_{ij}=\begin{cases}1/m_a,&i\in C_a,\ j\in C_{\pi(a)},\\
 0,&\text{otherwise}.
 \end{cases}
\end{equation}
Let $\Fset$ be the finite collection of matrices obtainable this way.
This definition uses a \emph{common partition of row and column labels};
arbitrary independent row and column relabellings do not describe the
same condition for matrix powers.

Associated with this partition is the orthogonal averaging projection
\begin{equation}\label{eq:Edef}
 E_{ij}=\begin{cases}1/m_a,&i,j\in C_a,\\0,&\text{otherwise}.
 \end{cases}
\end{equation}
It satisfies $E^2=E=E^{\mathsf T}$ and $EF=FE=F$. Also $F^k$ has the
form \eqref{eq:Fdef} with $\pi$ replaced by $\pi^k$. Consequently,
\begin{equation}\label{eq:pdef}
 p:=f(F)=f(E)=\prod_{a=1}^s\frac{m_a!}{m_a^{m_a}}>0,
 \qquad f(F^k)=p\quad(k\ge1).
\end{equation}
These matrices include permutations, $J_n$, direct sums of uniform
blocks, and periodic cycles of uniform blocks; they are not all
idempotent. Here $J_m$ means the $m\times m$ matrix with every entry
$1/m$.

\begin{theorem}\label{thm:main}
For every positive integer $n$, there is $N(n)\ge2$ such that for every
$A\in\Om$ and every integer $k\ge N(n)$,
\[
 f(A^k)\le f(A).
\]
Equality holds if and only if $A\in\Fset$.
\end{theorem}

The theorem compares every sufficiently large power with the
\emph{original matrix}, not necessarily with the immediately preceding
power. It is not a claim that $f(A^k)$ is monotone in $k$. The proof gives
an explicit local cutoff $16n^2$, but does \emph{not} establish that this
is a valid global cutoff.

\section{A row-averaging bound with a monotone defect}
Throughout the proof $\fn{\cdot}$ is the Frobenius norm and
$\op{\cdot}$ the Euclidean operator norm. Every $A\in\Om$ is a
Euclidean contraction: Jensen's inequality and the column sums give
\[
 \sum_i\left|\sum_j a_{ij}x_j\right|^2
 \le\sum_{i,j}a_{ij}|x_j|^2=\sum_j|x_j|^2.
\]
The same statement holds for row vectors by transposition.

Fix a partition and its projection $E$. For a doubly stochastic matrix
$M$ with $EM=M$, define
\begin{equation}\label{eq:variance}
 V_E(M):=n-\sum_{a=1}^s\sum_{j=1}^n
 \left(\sum_{i\in C_a}M_{ij}\right)^2.
\end{equation}
\begin{lemma}[Categorical occupancy]\label{lem:occupancy}
For every $M\in\Om$ satisfying $EM=M$,
\begin{equation}\label{eq:occupancybound}
 f(M)\le p\left(1-\frac{V_E(M)}{2n^2}\right)\le p,
 \qquad p=\prod_a\frac{m_a!}{m_a^{m_a}}.
\end{equation}
For a fixed $A\in\Om$, the sequence $V_E(EA^k)$ is nondecreasing.
\end{lemma}
\begin{proof}
Rows of $M$ agree within each class. Put
$y_{aj}=\sum_{i\in C_a}M_{ij}$, so $M_{ij}=y_{aj}/m_a$ for $i\in C_a$.
For each column $j$, independently choose a class $a$ with probability
$y_{aj}$. These are probabilities because $y_{aj}\ge0$ and
$\sum_a y_{aj}=1$. Let $N_a$ be the number of columns assigned to class
$a$. Its expectation is $\sum_j y_{aj}=m_a$.

An assignment of columns with exactly $m_a$ columns in each class has
$\prod_a m_a!$ row bijections contributing to the permanent. All such
bijections for this assignment have the same product of matrix entries.
Expanding the permanent therefore gives the exact identity
\begin{equation}\label{eq:occupancyidentity}
 \frac{f(M)}{p}=\mathbb P\{N_a=m_a\text{ for every }a\}.
\end{equation}
Independence across columns implies
\[
 \mathbb E\|N-m\|_2^2=\sum_{a,j}y_{aj}(1-y_{aj})=V_E(M).
\]
Since both $N$ and $m$ have nonnegative coordinates summing to $n$,
$\|N-m\|_2^2\le2n^2$. It vanishes on the event in
\eqref{eq:occupancyidentity}. Thus its expectation is at most
$2n^2$ times the probability of the complementary event, proving
\eqref{eq:occupancybound}.

Let $u_a$ be the column indicator of $C_a$. Since $u_a^{\mathsf T}E=
u_a^{\mathsf T}$,
\[
 V_E(EA^k)=n-\sum_a\|u_a^{\mathsf T}A^k\|_2^2.
\]
Right multiplication by $A$ does not increase any of these squared
norms. This proves monotonicity.
\end{proof}

\section{Perturbations of a uniform block-permutation matrix}
Fix $F\in\Fset$ with partition, $E$, and $p$ as above. In this section
$n\ge2$. For $A\in\Om$, write
\begin{equation}\label{eq:epsdelta}
 X=A-F,\qquad \eps=\fn X,\qquad
 \delta=\sum_{F_{ij}=0}A_{ij}.
\end{equation}
The quantity $\delta$ is the total mass leaking outside the support of
$F$. Cauchy--Schwarz gives $0\le\delta\le n\eps$.

\begin{lemma}[Tangent and leakage components]\label{lem:split}
There is an orthogonal decomposition $X=Y+Z$ such that $Y$ is supported
on $C_a\times C_{\pi(a)}$, has zero row and column sums in each such
block, and
\begin{equation}\label{eq:splitbounds}
 EY=YE=0,\qquad
 \eps^2=\fn Y^2+\fn Z^2,\qquad
 \fn Z\le\sum_{i,j}|Z_{ij}|\le4\delta.
\end{equation}
If $\eps\le1/n$, then $F+Y\in\Om$.
\end{lemma}
\begin{proof}
Orthogonally project each allowed block of $X$ onto the subspace with
zero row and column sums, and put $Y=0$ elsewhere. In an allowed
$m\times m$ block with row sums $r_i$, column sums $c_j$, and total $t$,
\[
 Y_{ij}=X_{ij}-\frac{r_i}{m}-\frac{c_j}{m}+\frac{t}{m^2}.
\]
This is the orthogonal projection, so the norm identity holds.
The zero block sums give $EY=YE=0$.

Every $r_i,c_j$ is nonpositive: it is minus the mass of $A$ outside the
allowed block in the corresponding row or column. In each block,
\[
 \sum_{i,j}|Z_{ij}|
 \le\sum_i|r_i|+\sum_j|c_j|+|t|=3(-t).
\]
Summing $-t$ over allowed blocks gives $\delta$. The entries of $Z=X$
outside the allowed blocks are nonnegative and total $\delta$, proving
the bound $4\delta$. Finally $\fn Y\le\eps$, so $\eps\le1/n$ keeps
all allowed entries of $F+Y$ nonnegative; its row and column sums remain
one.
\end{proof}

\subsection{Permanent derivatives and a lower bound for the starting matrix}
Set
\begin{equation}\label{eq:C}
 C=C_n:=\frac{n^2(n-2)!}{2}.
\end{equation}
On $\Om$, each second partial derivative of $f$ has absolute value at
most $(n-2)!$: it is either zero or the permanent of a submatrix with
entries at most one. Since $\sum|H_{ij}|\le n\fn H$, the Hessian
operator norm is at most $2C$. In particular, for a segment in $\Om$,
\begin{equation}\label{eq:taylor2}
 |f(M+D)-f(M)-Df_M(D)|\le C\fn D^2.
\end{equation}
For $n\ge3$ the third Taylor remainder in a direction $Y$ on such a
segment is bounded by
\[
 \frac{n^3(n-3)!}{6}\fn Y^3\le C\fn Y^3.
\]
For $n=2$ the third remainder is zero.

The permanental cofactor at an allowed position of $F$ equals $p$,
whereas at a position outside its support it is zero. Therefore
\begin{equation}\label{eq:linearper}
 Df_F(X)=-p\delta,\qquad Df_F(Y)=0,\qquad Df_F(Z)=-p\delta.
\end{equation}
For a tangent block $Y_a$ of size $m_a>1$, direct expansion at $J_{m_a}$
gives the quadratic coefficient
\begin{equation}\label{eq:hessianper}
 f(F+tY)=p+t^2p\sum_{a:m_a>1}
 \frac{m_a}{2(m_a-1)}\fn{Y_a}^2+O(t^3).
\end{equation}
To verify the coefficient, the sum of $Y_{ij}Y_{h\ell}$ over distinct
rows and columns equals $\fn{Y_a}^2$, because the row and column sums
vanish. Passing from ordered to unordered row pairs introduces $1/2$;
the complementary permanent equals $(m_a-2)!/m_a^{m_a-2}$.
Dividing by $m_a!/m_a^{m_a}$ gives the displayed factor. Blocks of size
one have no tangent directions, and cross-block linear products vanish.

If $\eps\le1/n$ and $C\eps\le p/4$, the remainder bound shows
\begin{equation}\label{eq:tangentlower}
 f(F+Y)\ge p+\frac p4\fn Y^2.
\end{equation}
The segment $F+Y+tZ$, $0\le t\le1$, lies in $\Om$, and its distance
from $F$ is at most $\eps$ by orthogonality. The Hessian bound and
\eqref{eq:splitbounds} therefore give
\[
 |f(A)-f(F+Y)-Df_F(Z)|\le2C\eps\fn Z\le8C\eps\delta.
\]
If also $8C\eps\le p$ and $\delta\le1/4$, then
\begin{align}
 f(A)&\ge p-2p\delta+\frac p4\fn Y^2\notag\\
 &\ge p-2p\delta+\frac p4\eps^2-4p\delta^2\notag\\
 &\ge p-3p\delta+\frac p4\eps^2.\label{eq:startinglower}
\end{align}

\section{An estimate uniform in arbitrarily large powers}
The difficulty is that $A^k$ need not stay near $F^k$ when $k$ is large.
The following argument differentiates a \emph{fixed} power only, then
uses a monotone quantity to cover all later powers.

\begin{lemma}[Fixed-time growth of the mixing defect]\label{lem:mixing}
For an integer $r\ge1$, set $g_r(A)=V_E(EA^r)$. Then
\begin{equation}\label{eq:gderivative}
 g_r(F+Y)=0,\qquad Dg_r(F)[X]=2r\delta.
\end{equation}
If $\eps\le1/n$ and $\eps\le1/(16nr)$, then
\begin{equation}\label{eq:growth}
 V_E(EA^k)\ge r\delta\qquad\text{for every }k\ge r.
\end{equation}
\end{lemma}
\begin{proof}
The class indicators satisfy $u_a^{\mathsf T}F=u_{\pi(a)}^{\mathsf T}$
and $u_a^{\mathsf T}Y=0$. Thus
$u_a^{\mathsf T}(F+Y)^r=u_{\pi^r(a)}^{\mathsf T}$, proving the first
identity in \eqref{eq:gderivative}.

The derivative of $A^r$ at $F$ is
\[
 \mathcal D_r(X)=\sum_{j=0}^{r-1}F^jXF^{r-1-j}.
\]
For each summand, multiplication by the powers of $F$ carries allowed
positions of $X$ to the support of $F^r$ and disallowed positions to its
complement. This follows by following the class permutation $\pi$;
the correspondence on classes is bijective. The multipliers are doubly
stochastic, so the total mass carried from disallowed positions is
exactly $\delta$. This remains valid when one of the powers is $F^0=I$.
The derivative therefore has outside mass $r\delta$ and inside total
$-r\delta$. Differentiating \eqref{eq:variance} at the deterministic
class assignment of $F^r$ gives $Dg_r(F)[X]=2r\delta$.

For the error estimate, let $U$ have rows $u_a^{\mathsf T}$; then
$\fn U^2=n$ and $g_r(A)=n-\fn{UA^r}^2$. On $\Om$, the operator-norm
bounds for the first and second derivatives of $A^r$ are respectively
$r\fn H$ and $r(r-1)\fn H\fn K$. The product rule yields
\[
 |D^2g_r(A)[H,K]|\le4nr^2\fn H\fn K.
\]
Comparing along $F+Y+tZ$ and using $g_r(F+Y)=0$ consequently gives
\begin{align*}
 |g_r(A)-Dg_r(F)[Z]|
 &\le4nr^2\eps\fn Z\le16nr^2\eps\delta.
\end{align*}
Also $Dg_r(F)[Y]=0$, so $Dg_r(F)[Z]=2r\delta$.
For $\eps\le1/(16nr)$ the error is at most $r\delta$.
This proves $g_r(A)\ge r\delta$. Lemma~\ref{lem:occupancy} makes the
same lower bound valid for every later power, proving \eqref{eq:growth}.
\end{proof}

\begin{lemma}[Uniform suppression of non-averaged rows]\label{lem:fast}
If $\eps\le1/2$, then for every integer $k\ge1$,
\begin{equation}\label{eq:fast}
 \fn{(I-E)A^k}\le8\delta+\eps^k.
\end{equation}
Moreover,
\begin{equation}\label{eq:rowtay}
 f(A^k)\le f(EA^k)+C\fn{(I-E)A^k}^2.
\end{equation}
\end{lemma}
\begin{proof}
Put $D_k=(I-E)A^k$, $R=(I-E)AE$, and $T=(I-E)A(I-E)$.
Since $(I-E)F=F(I-E)=0$ and $YE=0$,
\[
 \fn R=\fn{(I-E)ZE}\le4\delta,
 \qquad\op T\le\eps,
 \qquad\fn{D_1}\le\eps.
\]
The exact recurrence is $D_k=RA^{k-1}+TD_{k-1}$. Contraction of powers
of $A$ gives
\[
 \fn{D_k}\le4\delta\sum_{j=0}^{k-2}\eps^j+\eps^k
 \le8\delta+\eps^k.
\]
For $M=EA^k$, the rows within a class are equal, so their permanental
cofactor rows are equal as well. Since $ED_k=0$, the sum of the rows of
$D_k$ within each class is zero; hence $Df_M(D_k)=0$. Both $M$ and
$M+D_k=A^k$ belong to $\Om$. Applying \eqref{eq:taylor2} proves
\eqref{eq:rowtay}.
\end{proof}

\begin{theorem}[Explicit uniform local gap]\label{thm:local}
For $n\ge2$, $F\in\Fset$, and its associated $p$, set
\begin{equation}\label{eq:radius}
 r=16n^2,\qquad C=\frac{n^2(n-2)!}{2},\qquad
 \rho_F=\min\left\{\frac1{16nr},\frac{p}{128Cn}\right\}.
\end{equation}
For every $A\in\Om$ with $\eps=\fn{A-F}\le\rho_F$ and every integer
$k\ge r$,
\begin{equation}\label{eq:localgap}
 f(A)-f(A^k)\ge4p\delta+\frac p8\eps^2,
 \qquad \delta=\sum_{F_{ij}=0}A_{ij}.
\end{equation}
In particular the inequality is strict unless $A=F$.
\end{theorem}
\begin{proof}
The radius guarantees all hypotheses used in
\eqref{eq:startinglower}, including $\delta\le n\eps\le p/(128C)\le1/4$.
It also guarantees the hypotheses of Lemmas~\ref{lem:mixing} and
\ref{lem:fast}. Combining those lemmas with
\eqref{eq:occupancybound} gives, for every $k\ge r$,
\begin{align*}
 f(A^k)&\le p-\frac{pr\delta}{2n^2}
       +C(8\delta+\eps^k)^2\\
 &\le p-8p\delta+128C\delta^2+2C\eps^{2k}.
\end{align*}
The first error term is at most $p\delta$. Since
$\eps\le p/(128Cn)\le p/(16C)$, $\eps\le1$, and $2k-2\ge1$,
the second error term is at most $(p/8)\eps^2$. Therefore
\[
 f(A^k)\le p-7p\delta+\frac p8\eps^2.
\]
Subtract this from \eqref{eq:startinglower} to obtain
\eqref{eq:localgap}. No bound on $k\eps$ or $k\delta$ was imposed.
\end{proof}

\section{The global compactness argument}
We use the classical van der Waerden--Falikman--Egorychev theorem:
\begin{equation}\label{eq:vdw}
 f(T)\ge\frac{m!}{m^m}\quad(T\in\Omega_m),
 \qquad\text{with equality exactly when }T=J_m.
\end{equation}
The lower bound and its equality characterization are existing results,
not ingredients claimed as new here; see \cite{Gurvits} and its
references to Falikman and Egorychev.

\begin{lemma}[Limits of powers near a fixed matrix]\label{lem:global}
For any $A_*\in\Om$, there is a block-averaging projection $E$ as in
\eqref{eq:Edef} such that:
\begin{enumerate}[label=(\roman*)]
\item $f(A_*)\ge f(E)$, with equality exactly when $A_*\in\Fset$;
\item if $A_j\to A_*$, $k_j\to\infty$, and $A_j^{k_j}\to B$, with
$A_j\in\Om$, then $EB=B$ and consequently $f(B)\le f(E)$.
\end{enumerate}
\end{lemma}
\begin{proof}
Recall the standard finite Markov-chain cyclic decomposition, also used
for this problem in \cite{Mel}. A doubly stochastic matrix has no
transient states: an isolated sink strongly connected component in its
directed graph can have no incoming mass, by equality of total row and
column sums; remove such components successively. Thus, after a
\emph{simultaneous} row/column permutation, $A_*$ is a direct sum of
irreducible doubly stochastic matrices.

Within an irreducible component of period $d$, split into its cyclic
classes. Transitions go from each class to the next. The total outgoing
mass of one class equals its size, and the incoming mass of the next
equals that next class's size. These sizes are therefore all equal.
Each transition block $T_a$ between consecutive classes is square and
doubly stochastic. The usual convergence theorem for an irreducible
aperiodic finite chain, applied to the $d$-step restrictions, shows that
for a common multiple $L$ of the periods,
\begin{equation}\label{eq:Markovlimit}
 A_*^{Lt}\longrightarrow E\qquad(t\to\infty),
\end{equation}
where $E$ averages within every cyclic class. This is an orthogonal
projection because the stationary distribution is uniform.

Every perfect matching of the support must match each row class to
its unique succeeding column class. Hence the permanent factorizes as
$\prod_a f(T_a)$. Applying \eqref{eq:vdw} block by block gives
$f(A_*)\ge\prod_a m_a!/m_a^{m_a}=f(E)$. Equality forces every
$T_a=J_{m_a}$, exactly the definition of $\Fset$.

For the varying-matrix limit, fix an integer $t$. For all sufficiently
large $j$, $k_j\ge Lt$. Euclidean contraction implies
\[
 \op{(I-E)A_j^{k_j}}
 \le\op{(I-E)A_j^{Lt}}.
\]
Taking $j\to\infty$, then $t\to\infty$ in
\eqref{eq:Markovlimit}, gives $(I-E)B=0$. The limit $B$ is doubly
stochastic, so Lemma~\ref{lem:occupancy} gives $f(B)\le f(E)$.
There is no assumption that $k_j$ is divisible by $L$, and no assumption
of convergence uniform over the matrices $A_j$.
\end{proof}

\begin{proof}[Proof of Theorem~\ref{thm:main}]
For $n=1$ the assertion is immediate. Suppose for some fixed $n\ge2$
that no eventual threshold with the claimed strictness exists.
Then there are integers $k_j\to\infty$ and matrices
$A_j\in\Om\setminus\Fset$ such that
\[
 f(A_j^{k_j})\ge f(A_j).
\]
By compactness, pass to a subsequence with $A_j\to A_*$ and
$A_j^{k_j}\to B$. Lemma~\ref{lem:global} and continuity give
\[
 f(A_*)\le f(B)\le f(E)\le f(A_*).
\]
All inequalities are equalities. Thus $A_*=F\in\Fset$.
For large $j$, $\fn{A_j-F}\le\rho_F$ and $k_j\ge16n^2$.
Since $A_j\ne F$, Theorem~\ref{thm:local} then gives
$f(A_j^{k_j})<f(A_j)$, a contradiction. The exceptional matrices attain
equality for every positive exponent by \eqref{eq:pdef}.
\end{proof}

\section{A uniform quadratic stability consequence}
The same argument provides more than a strict inequality.
\begin{corollary}\label{cor:stability}
For every $n$ there are $c_n>0$ and $N_n\ge2$ such that for every
$A\in\Om$ and every integer $k\ge N_n$,
\begin{equation}\label{eq:globalstability}
 f(A)-f(A^k)\ge c_n\,\dist_{\mathrm F}(A,\Fset)^2.
\end{equation}
Neither the global $c_n$ nor $N_n$ is made explicit here.
\end{corollary}
\begin{proof}
If this were false for a fixed $n$, choose $k_j\ge j$ and $A_j\in\Om$
with
\[
 f(A_j)-f(A_j^{k_j})<\frac1j\dist_{\mathrm F}(A_j,\Fset)^2.
\]
No $A_j$ is exceptional. Pass to compact limits as before. The squared
distances are bounded, so the limiting permanent gap is nonpositive.
Lemma~\ref{lem:global} again forces the limit $A_*=F\in\Fset$.
For all sufficiently large $j$, Theorem~\ref{thm:local} gives
\[
 f(A_j)-f(A_j^{k_j})
 \ge\frac p8\fn{A_j-F}^2
 \ge\frac p8\dist_{\mathrm F}(A_j,\Fset)^2,
\]
contradicting $1/j<p/8$. For $n=1$ the distance vanishes and the
assertion is trivial.
\end{proof}

\section{Interpretation and scope}
The permanent of a transition matrix $A$ is the probability that $n$
independent walkers, started one at each of the $n$ states, occupy
\emph{distinct} states after one step. The corresponding probability
after $k$ steps is $f(A^k)$. Theorem~\ref{thm:main} says that, under
double stochasticity, every sufficiently late time has probability no
larger than the one-step probability, with a threshold depending only
on $n$.

The stronger assertion with $k=2$ is already false. Van Name's published
2022 example \cite{MO} is
\[
 A=\frac12\begin{pmatrix}
 1&0&1&0\\0&0&1&1\\1&1&0&0\\0&1&0&1
 \end{pmatrix},\qquad
 f(A)=\frac18<\frac9{64}=f(A^2).
\]
This example is included only as an attributed scope calibration, not
as a new counterexample. It is checked by direct permutation summation
and Ryser inclusion--exclusion in the package.

This is not a uniform mixing-time theorem. For example,
$A_\eta=(1-\eta)I+\eta J_n$ satisfies
$A_\eta^k=J_n+(1-\eta)^k(I-J_n)$, so convergence can be arbitrarily
slow as $\eta\downarrow0$. The monotone defect in
Lemma~\ref{lem:occupancy}, rather than a uniform mixing estimate, is
what makes the proof uniform.

The theorem applies to all real nonnegative doubly stochastic matrices,
including reducible, periodic, singular, irrational, and boundary
matrices. It does not extend the assertion to arbitrary row-stochastic,
complex, or signed matrices. It does not settle the distinct
Foregger--Sinkhorn tie-point problem, whose 2026 counterexample is due
to Lavi \cite{Lavi}, or the separate convex-mixture inequalities also
associated with Foregger's name.

\section{Exact diagnostics and review status}
The accompanying Python and C++ implementations independently test the
identities underlying the proof. Python uses subset dynamic programming
for permanents and dual-number binary exponentiation for power
derivatives. C++ uses Ryser inclusion--exclusion, a sequential recurrence
for those derivatives, explicit cofactor sums for the permanent Hessian,
and a block-matrix projection rather than row/column-mean subtraction.
The occupancy identity is checked by categorical count dynamic
programming in Python and direct assignment recursion in C++.

On the common 90 fixtures in dimensions two through seven, the two
implementations agree on 4,500 exact rational fields in 1,080 records.
They check 450 mixing derivatives, 90 tangent Hessians, 90 occupancy
identities and 450 specified powers. The fixture family is explicitly
defined; it is \emph{not} an exhaustive enumeration of doubly stochastic
matrices. Of the 90 fixtures, 69 are default fixtures that need not lie inside
the very conservative radius \eqref{eq:radius}; the other 21 are
radius-controlled fixtures. Their 42 late-power checks satisfy the explicit local gap, including exponents $16n^2+1$.
Seventeen deliberate mutations are rejected by the diagnostic guards. Finite checks neither prove the universal
quantifiers nor supply the unspecified global threshold.

All mathematical computations use integers or rational numbers. No
floating-point experiment, numerical spectral tolerance, or random
confidence statement is a premise of the proof. The initial fixture
selection is deterministic with a recorded seed; randomness only chooses
diagnostic examples.

The proof and code were produced and audited within one assistant
session. There has been no external specialist review, independent
institutional replication, or formal proof-assistant verification.
The dated primary-source search found no earlier publication of the
specific uniform theorem or proof, but cannot exclude unindexed or
private concurrent work. The bibliography and access limitations are
recorded separately in \texttt{sources/prior\_art.md}.

\begin{thebibliography}{9}
\bibitem{Minc} H. Minc, \emph{Permanents}, 1978, Conjecture 17.
Original attribution and catalogue location checked through
\cite{CW,Mel,Zhang}; the original page was not independently accessed.
\bibitem{CW} G.-S. Cheon and I. M. Wanless,
An update on Minc's survey of open problems involving permanents,
\emph{Linear Algebra and its Applications} \textbf{403} (2005), 314--342.
\href{https://doi.org/10.1016/j.laa.2005.02.030}{doi:10.1016/j.laa.2005.02.030}.
Author-hosted PDF inspected, especially p.~326, Conjecture 17.
\bibitem{Mel} K. Melnykova, \emph{Notes on Foregger's Conjecture},
M.Sc. thesis, University of Manitoba, 2012.
Official repository PDF inspected, especially Chapter 4 and
Theorem 4.10. Bibliographic year verified from the thesis front matter.
\bibitem{Chang90} D. K. Chang, On two permanental conjectures,
\emph{Linear and Multilinear Algebra} \textbf{26} (1990), 207--213.
Results attributed through \cite{CW,Mel,Zhang}; full original article
not independently inspected.
\bibitem{Zhang} F. Zhang, An update on a few permanent conjectures,
\emph{Special Matrices} \textbf{4} (2016), 305--316.
\href{https://doi.org/10.1515/spma-2016-0030}{doi:10.1515/spma-2016-0030};
\href{https://arxiv.org/abs/1608.02844}{arXiv:1608.02844}.
\bibitem{Gurvits} L. Gurvits, Van der Waerden/Schrijver--Valiant like
conjectures and stable (aka hyperbolic) homogeneous polynomials:
one theorem for all,
\href{https://arxiv.org/abs/0711.3496}{arXiv:0711.3496}, version 2.
Introduction states the van der Waerden bound and unique equality case,
crediting Falikman and Egorychev; an independent polynomial proof is given.
\bibitem{MO} J. Van Name (8 May 2022) and I. M. Wanless (15 May 2022),
answers to ``On permanent of a square of a doubly stochastic matrix,''
\href{https://mathoverflow.net/questions/422029/on-permanent-of-a-square-of-a-doubly-stochastic-matrix}{MathOverflow question 422029}.
Original, dated author discussion, not a peer-reviewed article.
Van Name supplies the $k=2$ counterexample; Wanless states the distinct
uniform-exponent problem as open in his account.
\bibitem{Lavi} Y. Lavi, A counterexample to the Foregger--Sinkhorn
tie-point conjecture, 13 August 2026,
\href{https://arxiv.org/abs/2608.13025}{arXiv:2608.13025}.
A different conjecture; not evidence that the powers problem was settled.
\end{thebibliography}
\end{document}
